\section{Doctrine and Theological Models}

Not every true theological sentence has the same authority. Several claims in
this study belong to the Church's received doctrinal synthesis:

\begin{itemize}
  \item Christ instituted the Eucharist at the Last Supper and commanded its
  celebration in his memory.
  \item By the consecration of bread and wine, their whole substances are
  changed into Christ's Body and Blood; the whole risen Christ is present
  truly, really, and substantially under either species.
  \item Calvary is Christ's unique, complete, once-for-all sacrifice.
  \item The Mass is a true sacrifice and the same sacrifice as Calvary: Christ
  is the same principal priest and victim, while the manner of offering is
  sacramental and unbloody.
  \item Eucharistic memorial makes the Paschal sacrifice present and applies
  its fruit; it is not bare mental recollection.
  \item The Church and faithful participate in and are joined to Christ's
  offering without supplementing its sufficiency.
  \item The Eucharist is also banquet, communion, ecclesial unity, and pledge
  of future glory.
\end{itemize}

Other formulations are theological models or explanatory syntheses. Saying
that the Supper ``anticipates'' Calvary is well grounded and used by modern
magisterial theology, but proposed accounts of \emph{how} a future sacrifice
can be sacramentally given require care. Saying that the sacrifice of Christ
is ``eternally present'' may intend his enduring heavenly priesthood or the
permanent efficacy of his completed offering; it becomes misleading if it
denies that the Passion occurred once in created history. Speaking of the
sacrifice of the Cross as made present in the Mass is sound when ``present''
is sacramental; it
becomes false if it implies renewed suffering or local replication of the
hill, crowd, and execution.

``The apostles were carried forward'' and ``we are carried back'' are still
more remote models. They are not dogma, conciliar language, or a conclusion
established by the Fathers. Their best content is already expressed more
securely by anticipation, anamnesis, representation, and participation.
Theological freedom does not require adopting an unnecessary mechanism,
especially where the mechanism generates new errors.

